%% GENERATED by bin/curate/props_ops_reference.py -- DO NOT EDIT BY HAND.
%% The op schemas in gui/schemas/operations/ are the source of truth; a runTests gate refuses a stale render.
\section{Operations reference: what the bench can be asked}\label{ch:props-ops}

Everything \code{choupoProps} does is one of the operations below, declared in \file{system/propsDict}:

\begin{lstlisting}[language=dict]
operations
(
    {
        name   myFirstOp;        // yours; names the output
        type   propertyScan1D;   // one of the types below
        ...                      // the type's own keys
    }
);
\end{lstlisting}

Operations run in the order written, each independent of the others.  In the tables, \textbf{*} marks a key the operation requires; everything else has a default or is optional.  Each entry names a runnable case, which is the fastest way to see the keys in a real dict.

\subsection{Evaluating a state}

\paragraph{\code{propertyPoint}}
Single-point property evaluation: at a given (T, P, composition) the thermo package reports every property the chosen models can compute (Psat, gamma\_i, K\_i, Cp\_ig, H, S, Z, v, ...). Results print to the run log.

\begin{tabular}{@{}p{0.30\linewidth}p{0.62\linewidth}@{}}
\hline
\code{properties} & Properties to evaluate \\
\code{state}\textbf{*} & Thermodynamic state \\
\hline
\end{tabular}


\begin{tutoriallink}
\file{tutorials/props/compare/acetone01\_ipa\_water\_azeotrope}
\quad (and 5 more)
\end{tutoriallink}

\paragraph{\code{propertyScan1D}}
One-dimensional property sweep: varies one state variable (T, P or a mole fraction) over n evenly-spaced points and evaluates the requested properties at each. Writes one CSV row per point.

\begin{tabular}{@{}p{0.30\linewidth}p{0.62\linewidth}@{}}
\hline
\code{output}\textbf{*} & Output file \\
\code{properties}\textbf{*} & Properties to evaluate \\
\code{state} & Held state \\
\code{vary}\textbf{*} & Scan axis \\
\hline
\end{tabular}


\begin{tutoriallink}
\file{tutorials/props/compare/acetone01\_ipa\_water\_azeotrope}
\quad (and 9 more)
\end{tutoriallink}

\paragraph{\code{propertyScan2D}}
Two-dimensional property sweep: varies two state variables on a regular grid and evaluates each property at every node. Writes one CSV row per node; failed evaluations are written as NaN.

\begin{tabular}{@{}p{0.30\linewidth}p{0.62\linewidth}@{}}
\hline
\code{output}\textbf{*} & Output file \\
\code{properties}\textbf{*} & Properties to evaluate \\
\code{state} & Held state \\
\code{varyX}\textbf{*} & X axis \\
\code{varyY}\textbf{*} & Y axis \\
\hline
\end{tabular}


\begin{tutoriallink}
\file{tutorials/props/scan/scan2d01\_co2\_compressibility}
\end{tutoriallink}

\paragraph{\code{propertyScanBinary}}
Binary liquid-liquid map at one (T, P): samples the mixing Gibbs energy g\_mix across the whole composition axis AND resolves the binodal by an LL flash of a declared feed. The two views teach each other — the g\_mix curve SHOWS the common tangent, the flash FINDS it. The feed must lie INSIDE the miscibility gap or the flash rightly reports one phase; for an asymmetric binodal that is not the equimolar point, which is exactly the lesson the water/1-butanol case carries.

\begin{tabular}{@{}p{0.30\linewidth}p{0.62\linewidth}@{}}
\hline
\code{feed} & LL-flash feed \\
\code{grid} & Grid resolution \\
\code{output} & CSV output \\
\code{state}\textbf{*} & Held state \\
\hline
\end{tabular}


\begin{tutoriallink}
\file{tutorials/props/scan/binary01\_lle\_water\_nbutanol}
\end{tutoriallink}

\paragraph{\code{propertyScanTernary}}
Ternary map at one (T, P): sweeps the composition simplex and reports the phase behaviour — the LLE solubility envelope with tie-lines, or a VLE/VLLE audit — from the declared activity model. The map is the model's fingerprint: change the declared model and the envelope moves.

\begin{tabular}{@{}p{0.30\linewidth}p{0.62\linewidth}@{}}
\hline
\code{grid} & Grid resolution \\
\code{mode} & Map mode \\
\code{output} & CSV output \\
\code{shard} & Simplex shard \\
\code{state}\textbf{*} & Held state \\
\code{tieStride} & Tie-line stride [-] \\
\hline
\end{tabular}


\begin{tutoriallink}
\file{tutorials/props/scan/ternary01\_audit\_vlle}
\quad (and 2 more)
\end{tutoriallink}

\paragraph{\code{purePhaseDiagram}}
Pure-component P-T phase diagram: the vapour-pressure curve from the component's own model, and — when the case declares the solid's data — the sublimation and melting curves meeting it at the triple point. The solid block is CASE data (axiom 4): triple point, heats of fusion and sublimation, and the fusion volume change (negative for water, which is why its melting line leans backwards — the lesson the reference case exists for).

\begin{tabular}{@{}p{0.30\linewidth}p{0.62\linewidth}@{}}
\hline
\code{grid} & Grid resolution \\
\code{output} & CSV output \\
\code{solid} & Solid-phase data \\
\hline
\end{tabular}


\begin{tutoriallink}
\file{tutorials/props/scan/phase01\_water\_full}
\end{tutoriallink}

\paragraph{\code{psychrometricChart}}
Psychrometric chart for ANY carrier + condensable pair, generated from the package's own models rather than copied from a handbook: saturation line, constant-relative-humidity curves and wet-bulb lines over a declared temperature span. Carrier and condensable are NOT interchangeable — only the condensable needs a vapour-pressure model, and the op refuses by name when it has none. Air-water is one instance, not the definition; the reference case runs N2-water.

\begin{tabular}{@{}p{0.30\linewidth}p{0.62\linewidth}@{}}
\hline
\code{P} & Chart pressure [Pa] \\
\code{TmaxC} & Upper temperature [-] \\
\code{TminC} & Lower temperature [-] \\
\code{carrier}\textbf{*} & Carrier (dry) gas \\
\code{condensable}\textbf{*} & Condensable component \\
\code{grid} & Grid resolution \\
\code{output} & CSV output \\
\code{relativeHumidity} & RH curves \\
\code{wetBulb} & Wet-bulb lines \\
\hline
\end{tabular}


\begin{tutoriallink}
\file{tutorials/props/heat/psychro01\_n2\_water}
\end{tutoriallink}

\paragraph{\code{steamTables}}
IAPWS-IF97 industrial water/steam properties. This is a propsDict operation under choupoProps (like propertyPoint), NOT a flowsheet unit. At startup it self-checks the transcribed coefficients against the release verification tables and REFUSES the run on any deviation above 1e-8. Give exactly ONE mode block per operation: point, saturation or isobar. CSV columns are canonical SI (h, u, s, cp, cv in J/kg-based units). Tutorial tutorials/props/steam/steam01\_if97\_verification locks every release verification table as a golden KPI at reltol 1e-8.

\begin{tabular}{@{}p{0.30\linewidth}p{0.62\linewidth}@{}}
\hline
\code{isobar} & Isobaric scan \\
\code{output} & CSV output \\
\code{point} & Single (p,T) state \\
\code{saturation} & Saturation-line scan \\
\hline
\end{tabular}


\begin{tutoriallink}
\file{tutorials/props/steam/steam01\_if97\_verification}
\end{tutoriallink}

\subsection{Mixtures and activity}

\paragraph{\code{activityCoefficients}}
Liquid activity coefficients of a MOLECULAR mixture at one temperature and composition, straight from the declared activity model. The glass-box counterpart of a VLE calculation: the gammas are what make a flash a flash, and this op prints them on their own so the model's behaviour can be read rather than inferred from a split. Change the declared model and the numbers move — which is the lesson.

\begin{tabular}{@{}p{0.30\linewidth}p{0.62\linewidth}@{}}
\hline
\code{composition}\textbf{*} & Liquid mole composition \\
\code{output} & CSV output \\
\code{temperature} & Temperature [K] \\
\hline
\end{tabular}


\begin{tutoriallink}
\file{tutorials/props/molecular/nrtl01\_ethanol\_water\_package}
\end{tutoriallink}

\paragraph{\code{gibbsMap}}
Equilibrium map over a (T, P) grid by Gibbs minimisation: at every grid point the declared species set equilibrates over the element balance, and the declared metric (a species' mole fraction, an element's recovery) is contoured. The chart's pedagogical weapons are the two DECLARED boxes: an \code{industrialWindow} (where industry actually operates — for ammonia, NOT where equilibrium is best, because the catalyst rules) and a \code{kineticBand} (the user-declared region where kinetics kill the approach), both watermarked with their labels so the map says why the plant is not at the thermodynamic optimum.

\begin{tabular}{@{}p{0.30\linewidth}p{0.62\linewidth}@{}}
\hline
\code{Pgrid}\textbf{*} & Pressure grid \\
\code{Tgrid}\textbf{*} & Temperature grid \\
\code{anchors} & Published anchor points \\
\code{elements}\textbf{*} & Elements \\
\code{feed}\textbf{*} & Feed moles \\
\code{industrialWindow} & Labelled (T, P) box \\
\code{kineticBand} & Labelled (T, P) box \\
\code{metric}\textbf{*} & Contoured metric \\
\code{output} & CSV output \\
\code{species}\textbf{*} & Species set \\
\code{temperatureApproach} & Equilibrium approach [K] \\
\hline
\end{tabular}


\begin{tutoriallink}
\file{tutorials/props/gibbs/nh3\_equilibrium\_map}
\end{tutoriallink}

\paragraph{\code{vleConsistency}}
Gibbs-Duhem thermodynamic-consistency test on MEASURED binary VLE data. Back-computes the activity coefficients from the data itself by modified Raoult (gamma\_i = y\_i P / (x\_i Psat\_i(T))), then reports the pointwise Gibbs-Duhem residual, the Herington area test (pass if |D-J| < 10) and the infinite-dilution gammas. Needs measured y: a bubble-T-vs-x set alone cannot run it, and testing an NRTL or Wilson CURVE is a tautology since those are consistent by construction. The CSV carries x1, lnGamma1, lnGamma2, lnRatio, gdResidual.

\begin{tabular}{@{}p{0.30\linewidth}p{0.62\linewidth}@{}}
\hline
\code{P} & Pressure [Pa] \\
\code{component}\textbf{*} & Reference component \\
\code{dataset}\textbf{*} & Measured dataset \\
\code{output} & CSV output \\
\code{partner} & Binary partner \\
\code{state} & Held state \\
\hline
\end{tabular}


\begin{tutoriallink}
\file{tutorials/props/compare/compare\_vle\_etoh\_water}
\end{tutoriallink}

\paragraph{\code{hConsistency}}
Enthalpy-surface identity audit: checks that the assembled package's enthalpy route is internally consistent for the named components at one temperature. A pure numerical identity on the package's own H surface — it needs no experimental data and answers a question the tutorials cannot: whether the rungs a component's record pins actually compose.

\begin{tabular}{@{}p{0.30\linewidth}p{0.62\linewidth}@{}}
\hline
\code{components}\textbf{*} & Components to audit \\
\code{output} & CSV output \\
\code{temperature} & Audit temperature [K] \\
\hline
\end{tabular}


\begin{tutoriallink}
\file{tutorials/props/thermoTest/h\_surface\_identity}
\end{tutoriallink}

\subsection{Electrolytes}

\paragraph{\code{speciate}}
Aqueous speciation: distributes the declared analytical totals over the curated equilibrium network and reports every species, its activity coefficient, the ionic strength, the water activity and the mineral saturation indices. By default it is SI-ONLY — it says which solids are supersaturated without precipitating any. Adding an \code{equilibrate \{ minerals ( ... ) \}} block lets the NAMED solids precipitate to their SI = 0 ceiling; naming the allowed set is the author's decision and is required, because which solid a brine is permitted to drop is chemistry, not a default.

\begin{tabular}{@{}p{0.30\linewidth}p{0.62\linewidth}@{}}
\hline
\code{activityModel} & Aqueous activity model \\
\code{analyticalTotals} & Analytical totals \\
\code{atmosphere} & Open-system gas partial pressures \\
\code{composition} & Composition form \\
\code{diagMeanIonic} & Cation/anion pairs to publish a mean ionic activity coefficient for \\
\code{diagSpecies} & Species to pin as diagnostics \\
\code{equilibrate} & Allowed precipitation \\
\code{network} & Admitted chemistry records \\
\code{networkScope} & Speciation network scope \\
\code{output} & CSV output \\
\code{pH}\textbf{*} & Given pH [-] \\
\code{reason} & Restriction reason \\
\code{temperature} & Temperature [K] \\
\code{totals} & Analytical totals (synonym) \\
\code{totalsBasis} & What the totals are \\
\code{vapour} & Vapour side (Edwards Eq 5/6/11) \\
\code{verifyGlobal} & Run the full catalogue self-check \\
\hline
\end{tabular}


\begin{tutoriallink}
\file{tutorials/props/electrolyte/edwards01\_sour\_water\_activity}
\quad (and 11 more)
\end{tutoriallink}

\paragraph{\code{pitzerActivity}}
Mean-ionic activity coefficient and osmotic coefficient of ONE salt by the Pitzer ion-interaction model, as a function of molality. The salt is named by its two ions and its stoichiometry is DERIVED from electroneutrality (Ca + Cl gives CaCl2), so the case never restates what the charges already say. The pair parameters come from the curated catalogue; a missing pair is an error, not a zero.

\begin{tabular}{@{}p{0.30\linewidth}p{0.62\linewidth}@{}}
\hline
\code{anion}\textbf{*} & Anion species \\
\code{cation}\textbf{*} & Cation species \\
\code{molality} & Molality scan \\
\code{output} & CSV output \\
\code{temperature} & Temperature [K] \\
\code{validation} & Reference points \\
\hline
\end{tabular}


\begin{tutoriallink}
\file{tutorials/props/electrolyte/archer01\_nacl\_cold\_to\_hot}
\quad (and 6 more)
\end{tutoriallink}

\paragraph{\code{electrolyteActivity}}
Mean-ionic activity and osmotic coefficients of the salt the CASE's declared thermophysical system carries — the package's own electrolyte surface, exercised directly. Where \code{pitzerActivity} names its salt by two ions and is the model's bench, this op takes the salt from the declaration and is the PACKAGE's bench: the numbers here are exactly what the flowsheet units will see, which is the point of running it before them.

\begin{tabular}{@{}p{0.30\linewidth}p{0.62\linewidth}@{}}
\hline
\code{molality} & Molality scan \\
\code{output} & CSV output \\
\code{temperature} & Temperature [K] \\
\hline
\end{tabular}


\begin{tutoriallink}
\file{tutorials/props/electrolyte/pitzer02\_nacl\_package}
\end{tutoriallink}

\paragraph{\code{enrtlMultiSalt}}
Multi-salt symmetric eNRTL bench: the mean-ionic gammas of two (or more) salts sharing an ion, swept across the mixing fraction at constant total molality — the Harned-rule experiment. Two formulations run side by side by declaration: \code{published} (the industrial expressions exactly as printed) and \code{refined} (gamma as the EXACT Gibbs-Duhem-consistent derivative of the same excess Gibbs energy, the chain rule kept through the composition-dependent mixing rules). The difference between them is measured, not asserted.

\begin{tabular}{@{}p{0.30\linewidth}p{0.62\linewidth}@{}}
\hline
\code{formulation} & Gamma formulation \\
\code{output} & CSV output \\
\code{salts}\textbf{*} & Salts \\
\code{sweepPoints} & Sweep points [-] \\
\code{tauEE} & Ion-ion tau [-] \\
\code{totalMolality} & Total molality [mol/kg] \\
\hline
\end{tabular}


\begin{tutoriallink}
\file{tutorials/props/electrolyte/enrtl\_multisalt\_nacl\_kcl}
\end{tutoriallink}

\paragraph{\code{enrtlMixedSolvent}}
Mixed-solvent eNRTL bench: the mean-ionic activity coefficient of a salt in water-alcohol mixtures, regressed system by system against published data. The case SPELLS OUT the binary tau parameters and the alcohol's own constants so it reads as its own citation — defaults exist (the published water-ethanol-NaCl set) but a case that states them teaches where every number came from. Each entry of \code{systems} is one solvent composition with its Debye-Hückel A\_phi and its measured (m, gamma) points; the op reports the deviation per system.

\begin{tabular}{@{}p{0.30\linewidth}p{0.62\linewidth}@{}}
\hline
\code{alcohol} & Alcohol constants \\
\code{output} & CSV output \\
\code{salt} & Salt tau parameters \\
\code{systems}\textbf{*} & Solvent compositions with data \\
\hline
\end{tabular}


\begin{tutoriallink}
\file{tutorials/props/electrolyte/enrtl\_mixed\_nacl\_ethanol\_esteso}
\end{tutoriallink}

\paragraph{\code{scalingScan}}
Concentration scan of a water analysis, the RO/evaporator scaling audit: the declared analysis is concentrated point by point along a recovery axis, the full speciation is solved at each point, and every mineral's saturation index is reported — WHERE each scale becomes possible, before any membrane is sized. Takes the same solution-definition keys as \code{speciate} (totals, pH incl. \code{solve}, atmosphere for an open system, temperature, equilibrate for allowed precipitation); adds the recovery axis and, with an equilibrate block, a \code{feedFlow} that turns precipitation into kg/day. \code{feedFlow} without \code{equilibrate} is REFUSED — it only feeds a rate an SI-only scan does not compute — and a bare number is refused too: it needs a volumetric unit.

\begin{tabular}{@{}p{0.30\linewidth}p{0.62\linewidth}@{}}
\hline
\code{activityModel} & Aqueous activity model \\
\code{analyticalTotals} & Water analysis \\
\code{atmosphere} & Open-system gas pins \\
\code{diagSpecies} & Species to pin as diagnostics \\
\code{equilibrate} & Allowed precipitation \\
\code{feedFlow} & Feed flow [m3/s] \\
\code{output} & CSV output \\
\code{pH} & pH \\
\code{recovery} & Recovery axis \\
\code{temperature} & Temperature [K] \\
\code{totals} & Water analysis (synonym) \\
\hline
\end{tabular}


\begin{tutoriallink}
\file{tutorials/props/electrolyte/pitzer\_vs\_davies\_ro\_concentrate}
\quad (and 4 more)
\end{tutoriallink}

\paragraph{\code{exchange}}
Ion-exchange equilibrium on the props bench: the declared water contacts a resin at a declared dose and the op reports the exchanged composition, the resin loading and the pH. The \code{exchange \{\}} block accepts resin, resinDose and CEC and NOTHING else — the capacity nameplate lives in the resin's own record, and an unknown key refuses by name with the grammar. \code{pH solve;} asks for the proton balance to be solved rather than given.

\begin{tabular}{@{}p{0.30\linewidth}p{0.62\linewidth}@{}}
\hline
\code{analyticalTotals} & Analytical totals \\
\code{diagSpecies} & Species to pin as diagnostics \\
\code{equilibrate} & Allowed precipitation \\
\code{exchange}\textbf{*} & Resin contact \\
\code{output} & CSV output \\
\code{pH} & pH \\
\code{temperature} & Temperature [K] \\
\code{totals} & Analytical totals (synonym) \\
\hline
\end{tabular}


\begin{tutoriallink}
\file{tutorials/props/electrolyte/softener\_brackish}
\end{tutoriallink}

\subsection{Fitting and estimating}

\paragraph{\code{fitParameters}}
Levenberg-Marquardt regression of thermophysical-model parameters against experimental data, and the SCORING of a package as it stands. Two modes, and the difference matters: \code{fit} (default) adjusts the listed parameters and writes a fit log (chi\textasciicircum{}2 and the parameter trajectory per iteration) plus a parity CSV at the optimum; \code{evaluate} changes nothing and only scores — 'how good are the pairs I already have against MY data?' — so in that mode a \code{parameters} list is OPTIONAL and, when present, is a pinned what-if whose entries carry \code{value} rather than the fit vocabulary of initial/min/max.

\begin{tabular}{@{}p{0.30\linewidth}p{0.62\linewidth}@{}}
\hline
\code{acceptance} &  \\
\code{dataset} &  \\
\code{evidence} &  \\
\code{mode} & Mode \\
\code{options} & Solver options \\
\code{output} & Output files \\
\code{parameters} & Parameters to vary \\
\code{residual}\textbf{*} & What is being fitted \\
\hline
\end{tabular}


\begin{tutoriallink}
\file{tutorials/props/adsorption/isotherm02\_fit\_jaschik}
\quad (and 7 more)
\end{tutoriallink}

\paragraph{\code{vaporPressureFit}}
Least-squares fit of measured vapour-pressure data to the Antoine form, producing the A, B and C a component record carries. Like \code{heatCapacityFit}, this is curation made visible: the residuals are reported beside the coefficients so the fit is judged, not just taken.

\begin{tabular}{@{}p{0.30\linewidth}p{0.62\linewidth}@{}}
\hline
\code{acceptance} &  \\
\code{component} & Component \\
\code{dataset} & Measured dataset \\
\code{evidence} &  \\
\code{output} & CSV output \\
\hline
\end{tabular}


\begin{tutoriallink}
\file{tutorials/props/curation/curate01\_water\_evidence\_partition}
\quad (and 1 more)
\end{tutoriallink}

\paragraph{\code{heatCapacityFit}}
Least-squares polynomial fit of measured heat-capacity data, producing the \code{coefficients ( ... )} list a component record carries. The fit is a CURATION act made visible: the residuals and the recovered coefficients are reported so the student can judge the fit before promoting it into a \code{.dat}.

\begin{tabular}{@{}p{0.30\linewidth}p{0.62\linewidth}@{}}
\hline
\code{acceptance} &  \\
\code{component} & Component \\
\code{dataset} & Measured dataset \\
\code{degree} & Polynomial degree [-] \\
\code{evidence} &  \\
\code{output} & CSV output \\
\hline
\end{tabular}


\begin{tutoriallink}
\file{tutorials/props/fit/cp01\_polynomial\_recovery}
\end{tutoriallink}

\paragraph{\code{estimateComponent}}
Group-contribution estimation of a component's constants — curation made visible, like the fit ops: the estimate is computed, compared against a reference where one is given, and written as a promotable draft, never silently into the catalogue.  THREE routes by \code{model}: a group method (Joback for small molecules, vanKrevelen for a polymer repeat unit) decomposes the molecule into \code{groups}; a scalar-anchored method (RiaziDaubert) builds a petroleum PSEUDO-component from \code{anchors \{ Tb; SG; \}} and needs no groups — the record it writes says loudly that a cut is a LUMP of hundreds of species, not a molecule; and the polymer route adds a \code{polymer \{\}} block whose packing factor k is visible, never hidden, because rho = M0/(k·Vw) and k IS the assumption.

\begin{tabular}{@{}p{0.30\linewidth}p{0.62\linewidth}@{}}
\hline
\code{anchors} & Anchor data \\
\code{component} & Component name \\
\code{derived} & Derived-property requests \\
\code{groups} & Structural groups \\
\code{model} & Estimation model \\
\code{output} & CSV output \\
\code{polymer} & Polymer options \\
\code{reference} & Reference values \\
\code{validation} & Reference values (synonym) \\
\hline
\end{tabular}


\begin{tutoriallink}
\file{tutorials/plant/polycaprolactonePlant/properties}
\quad (and 6 more)
\end{tutoriallink}

\subsection{Benches}

\paragraph{\code{kinetics1D}}
One-dimensional batch kinetics bench: integrates c(t) for an nth-order rate, overlays a measured dataset when one is declared, and — with a \code{fit} list — regresses the rate constant or the Arrhenius pair (k0, Ea) from the data. The MODEL-DISCRIMINATION lesson lives here: the same dataset run at order 1 and order 2 shows which residuals carry structure, which is adequacy, not consistency.

\begin{tabular}{@{}p{0.30\linewidth}p{0.62\linewidth}@{}}
\hline
\code{component} & Component \\
\code{dataset} & Measured dataset \\
\code{fit} & Parameters to fit \\
\code{output}\textbf{*} & CSV output \\
\code{rate}\textbf{*} & Rate law \\
\code{state} & State block \\
\code{time} & Time span \\
\hline
\end{tabular}


\begin{tutoriallink}
\file{tutorials/props/kinetics/rate01\_arrhenius\_recovery}
\end{tutoriallink}

\paragraph{\code{heatTransferBench}}
Self-check of the heat-transfer correlation library against its own textbook worked examples: every built-in correlation is evaluated at its published example's conditions and held to the printed answer within a declared tolerance. Runs with NO configuration — the bench IS the test — and the tolerances exist only to honour textbook rounding and chart-reading scatter, never to loosen a failing correlation.

\begin{tabular}{@{}p{0.30\linewidth}p{0.62\linewidth}@{}}
\hline
\code{aadTolerance} & Deviation tolerance [-] \\
\code{aadTolerancePhaseChange} & Phase-change tolerance [-] \\
\hline
\end{tabular}


\begin{tutoriallink}
\file{tutorials/props/heat/htbench01\_correlations}
\end{tutoriallink}

\paragraph{\code{isothermEval}}
Evaluates a curated adsorption isotherm — the loading of one adsorbate on one adsorbent as a function of pressure and temperature — from the adsorbent record's own parameters. The glass-box surface of what a fixed-bed or TSA unit uses internally, so the isotherm a breakthrough case runs on can be plotted and checked before the bed is ever solved.

\begin{tabular}{@{}p{0.30\linewidth}p{0.62\linewidth}@{}}
\hline
\code{adsorbate}\textbf{*} & Adsorbate \\
\code{adsorbent}\textbf{*} & Adsorbent \\
\code{grid}\textbf{*} & Evaluation grid \\
\code{output} & CSV output \\
\code{unitTwin} & Unit to mirror \\
\hline
\end{tabular}


\begin{tutoriallink}
\file{tutorials/props/adsorption/isotherm01\_eval\_13x\_co2}
\end{tutoriallink}

\paragraph{\code{elementalComposition}}
Glass-box surface of the ONE shared formula parser (src/thermo/ElementComposition) that every balance diagnostic uses. Decomposes chemical formulas into their per-element atom counts and reports them, so the parser a student relies on for the element balance can be exercised and read directly. Give \code{formulas}, \code{components}, or both; a formula it cannot parse is REFUSED by name rather than silently skipped.

\begin{tabular}{@{}p{0.30\linewidth}p{0.62\linewidth}@{}}
\hline
\code{components} & Catalogue components \\
\code{formulas} & Literal formulas \\
\code{output} & CSV output \\
\hline
\end{tabular}


\begin{tutoriallink}
\file{tutorials/props/thermoTest/atoms01\_formula\_parser}
\end{tutoriallink}

\subsection{Other}

\paragraph{\code{enthalpyConcentration}}
The saturated-liquid and saturated-vapour enthalpy loci of a BINARY at ONE pressure — the data an enthalpy-concentration (Ponchon-Savarit) construction is drawn on; the construction itself belongs to a drawing tool. Every row is one equilibrium: it carries BOTH ends of a tie line (the liquid x1 with h\_liquid, the vapour y1 with H\_vapour) at the ONE temperature they share, so a consumer never pairs rows by index and never interpolates one end of a tie. \code{role bubble} rows put x1 on the uniform grid, \code{role dew} rows put y1 on it; both are solved through the SAME bubble-point kernel, a dew node being the inversion y\_eq(x) = y1. The op REFUSES up front unless the thermophysical system has EXACTLY 2 components, and unless every component can be priced on the project's one enthalpy datum (the elements at 298.15 K): it names any component with no route to that datum, any component with no idealGasHeatCapacity\{\} block (no saturated-vapour leg), and any component whose standardThermochemistry.referenceState is off the ideal-gas rung — no private Cp integral is ever substituted. A node that cannot be solved (bubble Newton not converged, a non-monotonic y\_eq(x1) that has no single-valued dew inversion, a y1 outside the converged bubble image, an enthalpy kernel refusal) is PUBLISHED as a refusal row with its status, never dropped. Diagnostics headline lambda = H\_vapour − h\_liquid over the locus, the spread of which measures what constant molar overflow costs.

\begin{tabular}{@{}p{0.30\linewidth}p{0.62\linewidth}@{}}
\hline
\code{grid} & Composition grid \\
\code{output}\textbf{*} & CSV output \\
\code{state}\textbf{*} & The pressure of the diagram \\
\hline
\end{tabular}


\begin{tutoriallink}
\file{tutorials/props/scan/hxy01\_ethanol\_water\_1atm}
\end{tutoriallink}

\paragraph{\code{freezingPoint}}
Freezing-point curve T\_f(m) of a single salt in water, solved as chemical-potential equality between an actual crystallising SolidPhase (ice) and the Pitzer surface's water activity: f\_solid(T) = a\_w(m,T)*Psat(T). The first case-reachable consumer of the crystal; its announced approximations (no dCp term, sub-freezing Psat extrapolation) fire on every run. Diagnostics: T\_f and depression at the grid end, the dilute-limit slope against nu*K\_f (the derived cryoscopic constant), and the depression AAD against a measured dataset via the standard validation block.

\begin{tabular}{@{}p{0.30\linewidth}p{0.62\linewidth}@{}}
\hline
\code{anion}\textbf{*} & Anion species \\
\code{cation}\textbf{*} & Cation species \\
\code{molality} & Molality scan \\
\code{output} & CSV output \\
\code{temperature} & Temperature [K] \\
\code{validation} & Reference points \\
\hline
\end{tabular}


\begin{tutoriallink}
\file{tutorials/props/electrolyte/fpd01\_nacl\_freezing}
\end{tutoriallink}

\paragraph{\code{frictionBench}}
The Darcy friction factor's correlation library, in two parts. VERIFY (always): every registered correlation — Blasius, Colebrook-White, Haaland, Churchill — reproduces its own published anchor, and the deviation is printed beside the correlation's validity window and its primary citation. COMPARE (optional): all four are evaluated at the SAME Re and roughness, and the spread is published TWICE — over all of them, and over only those inside their stated windows. The two are reported apart because collapsing them would say 'the correlations disagree by 26 \%' when the truth is 'three agree within 2 \% and one was asked a question it was never fitted for'. The bench never says which correlation is right: that depends on the pipe, and ranking them would hide the choice the engineer has to make.

\begin{tabular}{@{}p{0.30\linewidth}p{0.62\linewidth}@{}}
\hline
\code{compare} & Comparison point \\
\code{deviationTolerance} & Anchor deviation tolerance [-] \\
\hline
\end{tabular}


\begin{tutoriallink}
\file{tutorials/props/hydraulics/moody01\_friction\_correlations}
\end{tutoriallink}

\paragraph{\code{phaseEnvelope}}
The P-T phase envelope of a MIXTURE at fixed feed composition: the bubble curve (V/F -> 0) and the dew curve (V/F -> 1) traced in the (T, P) plane and written as a long-form CSV (P\_Pa, T\_K, curve). An envelope is drawn to find the CRICONDENTHERM -- the highest temperature at which the feed can hold a liquid, what a pipeline is designed against -- and the retrograde region beside it, and both live at the nose, which no single-specification march can close: this trace does NOT close it, says so on every run, and names Michelsen's arc-length continuation with a switching specification (Fluid Phase Equilib. 4 (1980) 1) as the algorithm that would. The bubble and dew solves are BubblePoint::compute and DewPoint::compute -- the same routines the bubbleT and dewT units call, so neither saturation residual exists twice. Marching in pressure, the temperature maximum is not a turning point for the specification, so an INTERIOR maximum is reported as a BRACKETED cricondentherm naming the two pressures it lies between; an endpoint maximum is not reported, because it says where the range stops and not where the mixture peaks. A point that CONVERGED but is discontinuous with its branch is discarded and the branch ends there: the bubble curve is monotonically increasing in P up to the critical point (physics, not a heuristic) and the dew curve gets a continuity test, because both saturation solvers bracket T in 200-700 K and are warm-started, so Newton can land on a different root and report success -- which it did, drawing a bubble temperature FALLING with pressure and publishing a cricondentherm that was the resulting spike. A branch that produces no point at all is reported as never having started, which is a different finding from a branch that stopped, and has a different remedy.

\begin{tabular}{@{}p{0.30\linewidth}p{0.62\linewidth}@{}}
\hline
\code{composition}\textbf{*} & Feed composition \\
\code{output} & CSV output \\
\code{pressure}\textbf{*} & Pressure march \\
\hline
\end{tabular}


\begin{tutoriallink}
\file{tutorials/props/molecular/envelope01\_natural\_gas}
\quad (and 1 more)
\end{tutoriallink}

\paragraph{\code{reactionCurve}}
OPEN-LOOP IDENTIFICATION off a trajectory somebody else recorded. Given a CSV written by a choupoCtrl run, one column driven and one column watched, it cuts the record into step SEGMENTS, fits a first-order-plus-dead-time surrogate (K, tau, theta) to each by Nelder-Mead on the sum of squared residuals, and applies the classical tuning rules to the segment the reader nominates. It solves no plant and integrates nothing: the experiment must already exist on disk, and the op reads it. THE SURROGATE IS A SURROGATE — the record is whatever the plant did, an FOPDT is three numbers, and the statement that the second is standing in for the first reaches the reader at the site AND in the end-of-run caveat block. A segment whose response has not SETTLED yields no time constant, because both K and tau are read off the end of the curve; the steady-state test that decides this is declared, not assumed. A dead time smaller than the record's own sampling interval is UNRESOLVED, and every rule that divides by theta (Ziegler-Nichols, Cohen-Coon) is then REFUSED BY NAME for that segment — a controller gain computed from a number the experiment could not measure is a number with nothing behind it. IMC survives there, because its algebra never divides by theta. NOTHING IN CHOUPO EVER CLOSES A LOOP ON THESE SETTINGS: the published gains are what each rule says, not a claim that they work.

\begin{tabular}{@{}p{0.30\linewidth}p{0.62\linewidth}@{}}
\hline
\code{manipulated}\textbf{*} & The driven column \\
\code{measured}\textbf{*} & The watched column \\
\code{minimumPoints} & Shortest usable segment [-] \\
\code{output} & CSV output \\
\code{reference} & Which segment the headline quotes \\
\code{settling} & The steady-state test \\
\code{trajectory} & The recorded experiment \\
\code{tuning} & Tuning rules to apply \\
\hline
\end{tabular}


\begin{tutoriallink}
\file{tutorials/ctrl/ctrl17\_step\_reaction\_curve}
\end{tutoriallink}

\paragraph{\code{solubilityParameter}}
The Hildebrand solubility parameter, delta = sqrt((dHvap(T) - R T) / V\_m), for every component of the case, plus the pairwise |delta\_i - delta\_j| table that is what the question usually is: the smaller the gap, the more likely two liquids mix. Delta is DERIVED and never stored -- the latent heat comes from the same Watson correlation the enthalpy legs use, and the molar volume from the record's Vliq -- so a stored value would be a second home for a fact the record already fixes. Where a record declares a published delta as an ANCHOR, the deviation is printed beside the derived number; that column is a check on the derivation and never an input to it, and it exists because all three ways this arithmetic fails are silent (a dropped RT term, a wrong V\_m unit, and Pa\textasciicircum{}0.5 reported where the literature uses MPa\textasciicircum{}0.5 -- a factor of 31.6 that still reads as a solubility parameter). Refuses by name for a component with no Tc, no HvapTb or no Vliq, and at a temperature at or above Tc, where there is no liquid and delta is not defined. States on every call that Hildebrand carries no hydrogen bonding -- it cannot separate a polar solvent from a hydrogen-bonding one of the same cohesive energy density, so a close pair is a CANDIDATE and never a verdict -- and names Hansen's three-parameter split as what would, not implemented because its parameters are a dataset this project does not hold.

\begin{tabular}{@{}p{0.30\linewidth}p{0.62\linewidth}@{}}
\hline
\code{T} & Temperature [K] \\
\code{output} & CSV output \\
\hline
\end{tabular}


\begin{tutoriallink}
\file{tutorials/props/molecular/solubility01\_hildebrand\_ladder}
\end{tutoriallink}

\paragraph{\code{thielePellet}}
The PROPS BENCH for one catalyst pellet: it resolves the pellet from its asset record and its effective diffusivity from the ONE home that owns that resolution, solves the isothermal first-order intraparticle boundary-value problem numerically, and VERIFIES that answer against the closed form. It publishes the concentration field c(r)/c\_s, the Thiele modulus phi with the effectiveness factor eta it produces, the Weisz-Prater observable modulus eta*phi\textasciicircum{}2, and — swept — the classical eta(phi) family over the three geometries. The physics is NOT here: it lives in unitOperations/reactor/pellet/PelletDiffusion, which the four reactors that today assume eta = 1 will reach. The op judges nothing: it publishes phi, eta and the field, announces what it assumed (isothermal, first-order, one pellet size, external film neglected), and leaves the reader to decide whether diffusion matters in their case. A run whose numerical answer misses the declared verification tolerances exits 1.

\begin{tabular}{@{}p{0.30\linewidth}p{0.62\linewidth}@{}}
\hline
\code{P} & Pressure [Pa] \\
\code{T}\textbf{*} & Temperature [K] \\
\code{catalyst}\textbf{*} & Catalyst pellet \\
\code{diffusion}\textbf{*} & Diffusion route \\
\code{nodes} & Grid points [-] \\
\code{output} & CSV output \\
\code{rateConstant} & Rate constant, per pellet volume [1/s] \\
\code{rateConstantMass} & Rate constant, per catalyst mass [m3/(kg.s)] \\
\code{refinement} & Grid-refinement study \\
\code{species}\textbf{*} & Diffusing species \\
\code{sweep} & eta(phi) family \\
\code{thermal} & Temperature field (Prater) \\
\code{verification} & Verification tolerances \\
\hline
\end{tabular}


\begin{tutoriallink}
\file{tutorials/props/reactor/thiele01\_sphere\_first\_order}
\quad (and 1 more)
\end{tutoriallink}

